\chapter{Impl\'ementation du syst\`eme}

Ce chapitre pr\'esente l'impl\'ementation concr\`ete du syst\`eme Andon, en d\'ecrivant les \'etapes de d\'eveloppement, l'environnement de travail et la mise en \ oeuvre de chaque composant.

\section{Environnement de d\'eveloppement}

\subsection{Outils utilis\'es}

\begin{table}[htbp]
\centering
\renewcommand{\arraystretch}{1.3}
\begin{tabular}{l l}
\toprule
\textbf{Cat\'egorie} & \textbf{Outil} \\
\midrule
Syst\`eme d'exploitation & Windows 10/11 \\
Editeur de code & Visual Studio Code \\
Serveur backend & Node.js v18 LTS + npm \\
Base de donn\'ees & PostgreSQL 15 \\
Broker MQTT & Mosquitto 2.0 \\
Framework frontend & Angular 16 \\
Microcontr\^oleur & ESP32 (Arduino IDE) \\
Gestion de version & Git + GitHub \\
\bottomrule
\end{tabular}
\caption{Environnement de d\'eveloppement}
\label{tab:environnement_dev}
\end{table}

\subsection{Installation des d\'e-pendances}

L'installation du projet se fait en quelques commandes :

\begin{lstlisting}[caption={Installation du backend},label={lst:install_backend}]
# Cloner le depot
git clone https://github.com/sews/andon-iiot.git

# Installer les d\'ependances du serveur
cd server
npm install

# Installer les d\'ependances du client
cd ../client
npm install
\end{lstlisting}

\section{Mise en place de l'infrastructure}

\subsection{Installation de PostgreSQL}

La base de donn\'ees est cr\'e\'ee en ex\'e-cutant le script SQL d'initialisation :

\begin{lstlisting}[caption={Initialisation de la base de donn\'ees},label={lst:init_db}]
# Creation de la base
createdb sews_andon

# Execution du script de schema
psql -d sews_andon -f schema.sql

# Insertion des donnees de test
psql -d sews_andon -f seed.sql
\end{lstlisting}

\subsection{Configuration de Mosquitto}

Le broker MQTT est install\'e et configur\'e selon les besoins du syst\`eme :

\begin{lstlisting}[caption={Installation de Mosquitto},label={lst:install_mosquitto}]
# Installation (Ubuntu)
sudo apt install mosquitto

# Demarrage du service
sudo systemctl start mosquitto
sudo systemctl enable mosquitto

# Test de connexion
mosquitto_pub -t "test/topic" -m "hello"
\end{lstlisting}

\section{D\'eveloppement du firmware ESP32}

\subsection{Programme principal}

Le firmware de l'ESP32 est d\'evelopp\'e en Arduino (C/C++) et suit la logique suivante :

\begin{lstlisting}[caption={Boucle principale du firmware ESP32},
label={lst:firmware_main}]
void setup() {
  Serial.begin(115200);
  pinMode(BOUTON_PIN, INPUT);
  pinMode(LED_VERTE, OUTPUT);
  pinMode(LED_ROUGE, OUTPUT);
  
  connectWiFi();
  connectMQTT();
  
  digitalWrite(LED_VERTE, HIGH);
}

void loop() {
  if (!mqttClient.connected()) {
    reconnectMQTT();
  }
  mqttClient.loop();
  
  if (digitalRead(BOUTON_PIN) == HIGH) {
    delay(200); // Debounce
    envoyerAlerte();
    digitalWrite(LED_VERTE, LOW);
    digitalWrite(LED_ROUGE, HIGH);
  }
}
\end{lstlisting}

\subsection{Fonction d'envoi d'alerte}

\begin{lstlisting}[caption={Envoi d'une alerte Andon via MQTT},
label={lst:envoyer_alerte}]
void envoyerAlerte() {
  StaticJsonDocument<200> doc;
  doc["ligneId"] = LIGNE_ID;
  doc["timestamp"] =getTimestampTimestamp();
  doc["type"] = "panne";
  doc["urgence"] = "haute";
  
  char buffer[200];
  serializeJson(doc, buffer);
  
  String topic = "sews/andon/" 
    + String(LIGNE_ID) + "/alert";
  mqttClient.publish(
    topic.c_str(), buffer, true
  );
}
\end{lstlisting}

\section{D\'eveloppement du serveur backend}

\subsection{Structure du projet}

Le serveur backend suit la structure suivante :

\begin{lstlisting}[caption={Arborescence du serveur backend},
label={lst:server_structure}]
server/
  |-- server.js          Point d'entree
  |-- package.json        Dependances
  |-- config/
  |   |-- db.js          Configuration PostgreSQL
  |   |-- mqtt.js        Configuration MQTT
  |-- routes/
  |   |-- evenements.js  Routes API evenements
  |   |-- lignes.js      Routes API lignes
  |   |-- users.js       Routes API utilisateurs
  |-- middleware/
  |   |-- auth.js        Authentification JWT
  |-- models/
  |   |-- Evenement.js   Modele evenement
  |   |-- Ligne.js       Modele ligne
\end{lstlisting}

\subsection{D\'emarrage du serveur}

\begin{lstlisting}[caption={D\'emarrage du serveur backend},
label={lst:server_start}]
const express = require('express');
const http = require('http');
const { Server } = require('socket.io');
const mqtt = require('mqtt');
const { Pool } = require('pg');

const app = express();
const server = http.createServer(app);
const io = new Server(server);

// Connexion PostgreSQL
const pool = new Pool({
  host: 'localhost',
  port: 5432,
  database: 'sews_andon',
  user: 'sews_admin',
  password: process.env.DB_PASS
});

// Connexion MQTT
const mqttClient = mqtt.connect(
  'mqtt://localhost:1883'
);

server.listen(3000, () => {
  console.log('Serveur demarre sur :3000');
});
\end{lstlisting}

\section{D\'eveloppement du frontend Angular}

\subsection{Structure du projet}

\begin{lstlisting}[caption={Arborescence du frontend Angular},
label={lst:client_structure}]
client/
  |-- src/
  |   |-- app/
  |   |   |-- components/
  |   |   |   |-- dashboard/
  |   |   |   |-- lignes/
  |   |   |   |-- evenements/
  |   |   |-- services/
  |   |   |   |-- api.service.ts
  |   |   |   |-- socket.service.ts
  |   |   |-- models/
  |   |   |-- app-routing.module.ts
  |   |-- assets/
  |   |-- environments/
  |-- angular.json
  |-- package.json
\end{lstlisting}

\section{Tests et validation}

\subsection{Tests unitaires}

Des tests unitaires ont \'et\'e r\'ealis\'es pour chaque composant :

\begin{itemize}
    \item Tests de l'API REST (CRUD op\'erations).
    \item Tests de la connexion MQTT.
    \item Tests de la base de donn\'ees (requ\^etes SQL).
    \item Tests des composants Angular (navigation, affichage).
\end{itemize}

\subsection{Tests d'int\'e-gration}

Les tests d'int\'e-gration v\'erifient le bon fonctionnement de l'ensemble :

\begin{enumerate}
    \item Envoi d'un message MQTT de test.
    \item V\'erification de la cr\'eation de l'\'ev\'enement dans PostgreSQL.
    \item V\'erification de la diffusion via Socket.IO.
    \item V\'erification de l'affichage dans le frontend.
\end{enumerate}

\section{Synth\`ese du chapitre}

L'impl\'ementation du syst\`eme a mobilis\'e des comp\'etences dans diff\'erents domaines : programmation embarqu\'ee (Arduino/C++), d\'eveloppement backend (Node.js/Express), d\'eveloppement frontend (Angular/TypeScript) et administration de base de donn\'ees (PostgreSQL). Chaque composant a \'et\'e d\'evelopp\'e, test\'e et int\'egr\'e de mani\`ere coh\'erente.
